\documentclass[10pt,a4paper]{amsart}
\usepackage[margin=19mm]{geometry}
\usepackage{amsmath,amssymb,mathtools}
\usepackage[scr=rsfs,cal=euler]{mathalfa}
\usepackage{xfrac}
\usepackage[unicode,hidelinks,bookmarks=false]{hyperref}
\newcommand{\ie}{i.e.\ }
\newcommand{\cf}{cf.\ }
\newcommand{\resp}{respectively\ }
\newcommand{\Subsection}{\subsection}
\providecommand{\nobreakdash}{\nobreak\hbox{-}}
\setlength{\emergencystretch}{2em}
\multlinegap0pt
\usepackage{amssymb}
\usepackage{tikz}
\usepackage{float}
\newtheorem{theorem}{Theorem}
\newtheorem{lemma}{Lemma}
\newcommand{\gp}{{\rm gp}}
\newcommand{\lgp}{{\rm gp_{\rightarrow }}}
\title{The general position number of digraphs}
\author{Ullas Chandran S.V., Gabriele Di Stefano, Grahame Erskine, Haritha S, Elias John Thomas, James Tuite}
\date{Source: arXiv:2604.15909v1 (CC BY 4.0)}
\begin{document}
\maketitle

\noindent\emph{Source and licence.} The general position number of digraphs. Version arXiv:2604.15909v1 (CC BY 4.0), distributed by arXiv under the Creative Commons Attribution 4.0 International licence, http://creativecommons.org/licenses/by/4.0/, as recorded in the arXiv licence metadata for this exact version and shown on the abstract page https://arxiv.org/abs/2604.15909v1. Full licence notice: \texttt{licenses/arxiv-2604.15909-CC-BY-4.0.txt}.\par\smallskip

\noindent\textit{Editorial scope.} Counted unit: the article's theorem theorem (source0.tex lines 1179--1181). The article's complete original proof of it is delivered with it (source0.tex lines 1182--1193, one \texttt{proof} environment of 12 lines). No mathematics is written, repaired or re-derived: the statement and the proof are the article's own lines, in the article's own notation, and no retained line is substituted or re-flowed.  Retained uncounted as the source-local context the proof needs: lines 444--446; lines 993--995.  Nothing was omitted from any retained span, so the package carries no omission record.  No retained line contains a citation, so the wrapper carries no bibliography and no bibliography entry.  No retained line contains a figure environment, an image-inclusion command or drawing code, so no image is delivered and the package declares no asset dependency.  Editorial formatting only, in addition: an AMS \texttt{amsart} preamble that restates the article's own declarations and macros used by the retained lines, namely the environments \texttt{theorem}, \texttt{lemma} on separate counters, together with the standard packages those lines need.  The retained environments print in this document as Theorem 1, Lemma 1 in order of appearance, whereas the article prints the same items with the numbering of its own full document.  The complete native source of the article as distributed in the arXiv e-print for this exact version is bundled unchanged as \texttt{sources/198647/source0.tex}.
	\begin{theorem}\label{thm:isometric cover}
		If $\{ D_1,D_2,\dots,D_k\}$ is an isometric cover of a digraph $D$, then $\gp(D) \leq \sum_{i=1}^{k} \gp (D_i)$.
	\end{theorem}

	\begin{lemma}\label{lem:path cover bound}
		For any graph $G$, if $V(G)$ can be partitioned into $r$ paths of length at least 1 and $s$ copies of $K_1$, then $\lgp (G) \leq 2r+s$. In particular, $\lgp (G) \leq 2\rho (G)$.
	\end{lemma}

	\begin{theorem}
		The complete bipartite graph $K_{n,n}$ is $\gp $-full for $n \geq 2$.
	\end{theorem}

	\begin{proof}
		Write the vertex set of $K_{n,n}$ as $X \cup Y$, where $X = \{ x_1,\ldots ,x_n\} $ and $Y= \{ y_1,\ldots ,y_n\} $. For even $0 \leq r \leq n-2$, define the orientation $K^{(1)}(n)_r$ of $K_{n,n}$ as follows: the arcs directed from $Y$ to $X$ are:
		\begin{itemize}
			\item $y_i \rightarrow x_i$ for $1 \leq i \leq n$,
			\item $y_i \rightarrow x_j$ when $r+1 \leq j \leq i-2$, 
			\item $y_i \rightarrow x_j$ when $1 \leq i \leq r$, $i$ is odd and $j \geq i+2$, and
			\item $y_i \rightarrow x_j$ when $1 \leq i \leq r$, $i$ is even and $j \geq i+1$.
		\end{itemize}
		All other arcs are directed from $X$ to $Y$. Observe that for odd $i \leq r$, the subdigraph induced by $\{ x_i,x_{i+1},y_i,y_{i+1}\} $ is isometric, as is the subdigraph induced by $\{ x_{r+1},\ldots ,x_n\} \cup \{ y_{r+1},\ldots ,y_n\} $. Each of these subdigraphs has general position number 2 (observing that the last of these has the orientation defined by Lemma~\ref{lem:path cover bound}), so by Theorem~\ref{thm:isometric cover} we have $\gp (K^{(1)}(n)_r) \leq r+2$. The set $\{ x_1,\ldots ,x_r\} \cup \{ y_{r+1},y_{r+2}\} $ is in general position, so in fact $\gp (K^{(1)}(n)_r) = r+2$. This covers all even values in the range $[2,n]$. A similar idea works for the odd values. In the previous construction, when $0 \leq r \leq n-3$ is even, change the direction of the arc $x_{r+1} \rightarrow y_{r+2}$ to $y_{r+2} \rightarrow x_{r+1}$. Lemma~\ref{lem:path cover bound} and Theorem~\ref{thm:isometric cover} show that the general position number of this new orientation is at most $r+3$, and $\{ x_1,\ldots ,x_{r+2}\} \cup \{ y_{r+1}\} $ is a general position set of that order.
		
		For even $r$ in the range $0 \leq r \leq n-2$ we now derive an orientation $K^{(2)}(n)_r$ from $K^{(1)}(n)_r$ by orienting each arc between $\{ y_{r+1},\ldots ,y_n\} $ and $\{ x_{r+1},\ldots ,x_n\} $ towards $X$, i.e.\ $y_i \rightarrow x_j$ is an arc for $r+1 \leq i,j \leq n$. The subdigraph induced by $\{ x_{r+1},\ldots ,x_n\} \cup \{ y_{r+1},\ldots ,y_n\} $ is isometric and has general position number $2(n-r)$ (as each vertex is a source or sink), and so Theorem~\ref{thm:isometric cover} gives $\gp (K^{(2)}(n)_r) \leq 2n-r$, and $X \cup \{ y_{r+1},\ldots ,y_n\} $ is a general position set, so that equality holds. This gives all even values of the general position number in the range $[n+1,2n]$. As in the previous case, changing the arc $y_{r+1} \rightarrow x_{r+1}$ to $x_{r+1} \rightarrow y_{r+1}$ deals with the odd values in the range.   
	\end{proof}

\end{document}
